\documentclass{article}
\usepackage{amsmath,amssymb,hyperref}
\title{Retained comma presentation of variable-shape processes}
\date{15 November 2026}
\begin{document}\maketitle
\paragraph{Status.} Exact bridge to the original presented context/content theory, independently reviewed in the stated scope. The construction is not a claim that totalization preserves arbitrary pseudocoends. It separates a directly proved presentation recovery from a false independent-shape factorization.

\paragraph{Original data.} Fix the original pure interfaces $A,B$ and its literal native hom category
\[
 H(A,B)=(e_A:\mathcal M_A\to\mathcal C\ \downarrow\ d_B:\mathcal N_B\to\mathcal C).
\]
An object is $(P,Q,f:e_AP\to d_BQ)$; its arrow $(r,s)$ satisfies $d_B(s)f=f'e_A(r)$. Regard this category as a locally discrete bicategory. No faithfulness of $e_A$ or $d_B$ is assumed. When discussing composition and crossing, require the original square-functorial representative presentation, exactly as in the existing diagram-lifting theorem. This is not automatic from a multiplication on isomorphism classes.

\paragraph{Exact presented total.} Define objects by the data
\[
 (J,P:J\to\mathcal M_A,Q:J\to\mathcal N_B,
       f:e_AP\Rightarrow d_BQ).
\]
An arrow to $(I,P',Q',f')$ is a functor $u:I\to J$ and natural transformations $r:Pu\Rightarrow P'$, $s:Qu\Rightarrow Q'$ such that
\[
 d_B(s)\,f_u=f'e_A(r).
\tag{1}
\]
A two-cell from $(u,r,s)$ to $(v,r',s')$ is a natural transformation $\eta:u\Rightarrow v$ satisfying
\[
 r=r'(P\eta),\qquad s=s'(Q\eta).
\tag{2}
\]
These are equalities of actual grade transformations, not just their extension images. The body equality follows from (1), (2) and naturality of $f$. Composition uses the reverse-shape functor composite and the whiskered grade transformations. Two-cell compositions use naturality and composition in the grade categories.

There is an isomorphism between this explicitly presented two-category and the reverse-shape diagram two-category of ordinary functors $J\to H(A,B)$. Send an object to the diagram with value $(P_j,Q_j,f_j)$; naturality of $f$ is exactly its native body square on a $J$ arrow. An arrow becomes the natural transformation with components $(r_i,s_i)$. The modification condition in the locally discrete bicategory is precisely (2). These assignments are inverse on every displayed datum and commute with compositions. If unnormalized pseudodiagrams into a locally discrete bicategory are used instead, their invertible constraint two-cells force the ordinary functor equations, giving the same statement with the chosen unit frames. Thus the whole native context/content presentation, including potentially nonfaithful grade maps, is retained, not merely its observed bodies.

\paragraph{The valid factorization is shape-synchronized.} For a fixed $J$, the existing diagram presentation uses the pure base $\mathcal C^J$. At constant endpoint diagrams $A,B$ its comma category is exactly the identity-shape fibre of the presentation above: arrows have $u=\mathrm{id}_J$, and only the identity shape two-cell is retained. This is not the full sub-bicategory on all objects with shape $J$, which also has shape endofunctors. Its retained normalization is
\[
 M^J\odot_{\mathcal C^J}N^J\simeq H^J.
\tag{3}
\]
Here the middle object is an entire diagram, not a single object of $\mathcal C$. The normal representative of $(P,Q,f)$ has middle diagram $e_AP$, first body identity and second body $f$. A native $(r,s)$ is normalized by its content arrow, balance along the whole natural transformation $e_A r$, and its context arrow. Balance identity/composition give the inverse and triangle equations. This is the existing direct proof of (3), reused rather than replaced by a claim that exponentiation preserves coends.

Restriction along $u:I\to J$ sends this normal representative to the normal representative of $(Pu,Qu,f_u)$; balance along $e_A r$ restricts to the corresponding balance. Shape transformations use the retained grade maps $P\eta,Q\eta$, not quotients of them. Consequently one may totalize these already normalized fibre presentations with the data (1)--(2). This gives the presented total just proved. It is a shape-indexed normalization statement, not an ordinary relative tensor over the original fixed pure category. The diagram bases used in (3) are explicit mathematical bookkeeping for internal factorization; no new arbitrary internal shape map is thereby identified with an original external pure map.

\paragraph{Why independent product shapes do not supply that normalization.} If the context and content parts of a $J$-indexed object are separately assembled by the new product-shape operation, their shape is $J\times J$, whereas the correlated object has shape $J$. Recovery requires the diagonal restriction arrow. That restriction need not be an equivalence. Even with all body and grade categories terminal, take $J$ discrete with two objects. The independent diagram has four indexing objects, and the correlated diagram has two. An equivalence in the reverse-shape total two-category would induce an equivalence of these indexing categories, because its invertible unit/counit two-cells include natural isomorphisms of the shape composites. No such equivalence exists. Thus the proposed naive normalization fails already in the terminal-body case; no effect or observation issue is responsible.

This counterexample does not contradict (3): the latter deliberately uses a common shape and balances in $\mathcal C^J$. It also does not prove that no larger free completion could enforce a different factorization. It shows that the actual native product-shape construction does not make this diagonal invertible, and should not do so by silently losing synchronization data.

\paragraph{Crossing and composition recovery.} With the stipulated square-functorial original crossing, the existing diagram-lifting theorem constructs the common-shape crossing on (3) and proves its four comparisons and coherence on full representative generators. This supplies a retained synchronized presentation in every fibre. For independent product-shape global assembly, pull each input grade diagram back along its product projection and apply the original representative composition at each tuple. The original total-square functoriality provides the diagram arrows; its native/pure balance naturality supplies the corresponding transformations. Thus the result agrees with the already constructed product-shape assembly on the original process hom category, under the exact presentation isomorphism above. The equality is checked on native grade maps and bodies, not on body images alone.

The terminal shape recovers the original one-sided factors, mixed normalization and crossing, including their pure endpoint actions. Arbitrary external endpoint changes use their original total-grade square maps componentwise. This is stronger than retaining an unnamed hom category, but is not a new independent existence criterion for the crossing. It extends a crossing that has the required verified presentation.

\paragraph{Feedback boundary.} No new trace law follows merely from (3). A supplied native feedback operation that is functorial on the full presentation can be applied to diagrams; its diagram arrows are then part of the data, not inferred from endpoint equality. The separate fixed-state canonical theorem supplies one such operation and its synchronization preservation. A varying-state family requires its actual native transformation data and an additional preservation or directed-comparison argument; shape-synchronized normalization alone provides neither. In particular neither the terminal-body diagonal obstruction nor its repair by (3) settles the general instantaneous-feedback problem.

\paragraph{Attribution.} This is a direct presentation unpacking plus the already proved \href{https://github.com/eternal-november/eos-research-programmes/blob/main/projects/contextual-composition/foundations/diagram-native-interaction/Diagram_Lifting_of_Presented_Native_Interaction.tex}{diagram native interaction theorem}, whose normalization balances over $\mathcal C^J$ and whose shape restriction is established. The additional comparison is with the reverse-shape totalization and its product-shape assembly. No novelty claim is made for comma categories, diagram lifting, or the normalization mechanism.
\end{document}
